
\vspace{1cm}

\noindent  \textbf{Leia as instruções antes de fazer a lista:}

\vspace{0.3cm}
\noindent A resolução deve ser entregue em formato Jupyter Notebook ou Google Colab na plataforma Google Classroom, na tarefa designada e até o dia determinado.

\noindent  Não é necessário testar se os dados passados por argumento são válidos a menos que pedido na questão.


\vspace{0.5cm}


\begin{enumerate}[label=\textbf{Q\arabic*.}]

% ---------------------------------------------------------
\item Escreva uma função chamada \code{contar\_caracteres(texto)} que receba como parâmetro uma string qualquer. A função deve retornar um dicionário no qual as chaves sejam os caracteres individuais da string (convertidos para minúsculas e desconsiderando espaços em branco) e os valores correspondam à quantidade total de ocorrências de cada caractere. Teste a sua função de acordo com as entradas e saídas esperadas a seguir:

\begin{itemize}
    \item \textbf{Entrada:} \code{"A B C a b a"} $\rightarrow$ \textbf{Saída:} \code{\{'a': 3, 'b': 2, 'c': 1\}}
    \item \textbf{Entrada:} \code{"Marinha do Brasil"} $\rightarrow$ \textbf{Saída:} \code{\{'m': 1, 'a': 3, 'r': 2, 'i': 2,}\\\code{ 'n': 1, 'h': 1, 'd': 1, 'o': 1, 'b': 1, 's': 1, 'l': 1\}}
\end{itemize}


\ifmostrarGabarito{
\begin{minted}{python}
# Resposta Questão 1
def contar_caracteres(texto):
    frequencia = {}
    texto_limpo = texto.lower().replace(" ", "")
    for char in texto_limpo:
        frequencia[char] = frequencia.get(char, 0) + 1
    return frequencia
\end{minted}
}


% ---------------------------------------------------------
\item Considere um dicionário onde as chaves representam o nome de alunos e os valores representam suas respectivas turmas (exemplo: \code{\{"Ana": "Turma A", "Bruno": "Turma B",}\\\code{"Carla": "Turma A"\}}). Escreva uma função \code{inverter\_dicionario(d)} que receba esse dicionário e retorne um novo dicionário onde as chaves sejam o nome das turmas e os valores sejam listas contendo os nomes dos alunos matriculados em cada uma delas. Teste a sua função de acordo com as entradas e saídas esperadas a seguir:

\begin{itemize}
    \item \textbf{Entrada:} \code{\{"Ana": "Turma A", "Bruno": "Turma B", "Carla": "Turma A",}\\\code{"Daniel": "Turma B"\}}
    \item \textbf{Saída:} \code{\{'Turma A': ['Ana', 'Carla'], 'Turma B': ['Bruno', 'Daniel']\}}
\end{itemize}


\ifmostrarGabarito
\begin{minted}{python}
# Resposta Questão 2
def inverter_dicionario(d):
    resultado = {}
    for aluno, turma in d.items():
        if turma not in resultado:
            resultado[turma] = []
        resultado[turma].append(aluno)
    return resultado
\end{minted}
\fi



% ---------------------------------------------------------
\item Dado o dicionário a seguir, que representa um estoque de produtos e seus respectivos preços:
\begin{minted}{python}
estoque = {
    "camisa": 45.0,
    "calça": 120.0,
    "meia": 12.5,
    "casaco": 250.0,
    "boné": 35.0
}
\end{minted}
Escreva um trecho de código em Python utilizando laço de repetição e estrutura condicional para construir um novo dicionário denominado \code{produtos\_caros}, contendo apenas os itens que possuem preço estritamente superior a R\$ 50,00. Teste a sua função de acordo com as entradas e saídas esperadas a seguir:

\begin{itemize}
    \item \textbf{Entrada:} O dicionário \code{estoque} fornecido.
    \item \textbf{Saída:} \code{\{'calça': 120.0, 'casaco': 250.0\}}
\end{itemize}

\ifmostrarGabarito
\begin{minted}{python}
# Resposta Questão 3
estoque = {
    "camisa": 45.0,
    "calça": 120.0,
    "meia": 12.5,
    "casaco": 250.0,
    "boné": 35.0
}

produtos_caros = {}
for produto, preco in estoque.items():
    if preco > 50.0:
        produtos_caros[produto] = preco

print(produtos_caros)
\end{minted}
\fi


% ---------------------------------------------------------
\item Crie uma função \code{agrupar\_por\_tamanho(frase)} que receba uma frase (string). A função deve tratar a frase separando as palavras, ignorando maiúsculas e pontuações simples. O retorno deve ser um dicionário onde:
\begin{itemize}
    \item As chaves sejam números inteiros representando o comprimento (tamanho) das palavras.
    \item Os valores sejam listas com as palavras da frase que possuem aquele comprimento exato.
\end{itemize}

Teste a sua função de acordo com as entradas e saídas esperadas a seguir:

\begin{itemize}
    \item \textbf{Entrada:} \code{"O sol brilhou e o sol se foi."}
    \item \textbf{Saída:} \code{\{1: ['o', 'e', 'o'], 3: ['sol', 'sol', 'foi'], 7: ['brilhou'], }\\\code{2: ['se']\}}
\end{itemize}

\ifmostrarGabarito
\begin{minted}{python}
# Resposta Questão 4
import string

def agrupar_por_tamanho(frase):
    frase_limpa = frase.translate(str.maketrans('', '', string.punctuation)).lower()
    palavras = frase_limpa.split()
    
    grupos = {}
    for palavra in palavras:
        tamanho = len(palavra)
        if tamanho not in grupos:
            grupos[tamanho] = []
        grupos[tamanho].append(palavra)
        
    return grupos
\end{minted}
\fi



% ---------------------------------------------------------
\item Considere a estrutura de dicionário aninhado contendo as notas de estudantes em três disciplinas distintas:

\begin{minted}{python}
alunos = {
    "Ana": {"Matemática": 8.5, "Português": 9.0, "História": 7.5},
    "Bruno": {"Matemática": 5.0, "Português": 6.0, "História": 5.5},
    "Carla": {"Matemática": 9.5, "Português": 8.0, "História": 9.0}
}
\end{minted}

Escreva um programa que percorra a estrutura e gere um novo dicionário \code{relatorio\_final}. Para cada aluno, calcule a média aritmética das notas e atribua o status \code{"Aprovado"} se a média for maior ou igual a 7.0, ou \code{"Recuperação"} caso contrário.

Teste a sua função de acordo com as entradas e saídas esperadas a seguir:

\begin{itemize}
    \item \textbf{Entrada:} O dicionário \code{alunos} fornecido.
    \item \textbf{Saída:} \code{\{'Ana': \{'media': 8.33, 'status': 'Aprovado'\}, 'Bruno': }\\\code{\{'media': 5.5, 'status': 'Recuperação'\}, 'Carla': \{'media': 8.83, }\\\code{'status': 'Aprovado'\}\}}
\end{itemize}

\ifmostrarGabarito
\begin{minted}{python}
# Resposta Questão 5
alunos = {
    "Ana": {"Matemática": 8.5, "Português": 9.0, "História": 7.5},
    "Bruno": {"Matemática": 5.0, "Português": 6.0, "História": 5.5},
    "Carla": {"Matemática": 9.5, "Português": 8.0, "História": 9.0}
}

relatorio_final = {}
for aluno, disciplinas in alunos.items():
    notas = disciplinas.values()
    media = round(sum(notas) / len(notas), 2)
    status = "Aprovado" if media >= 7.0 else "Recuperação"
    relatorio_final[aluno] = {"media": media, "status": status}

print(relatorio_final)
\end{minted}
\fi



% ---------------------------------------------------------
\item Dada uma lista de vendas onde cada item é uma string no formato \code{"PRODUTO:QUANTIDADE"}:
\begin{minted}{python}
vendas = ["maçã:5", "banana:3", "maçã:2", "uva:1", "laranja:8", "banana:4"]
\end{minted}
Escreva um algoritmo em Python que processe os dados da lista utilizando divisão de strings (\code{.split()}) e acúmulo em um dicionário, resultando em uma contagem total consolidada da quantidade vendida de cada produto.

Teste a sua função de acordo com as entradas e saídas esperadas a seguir:

\begin{itemize}
    \item \textbf{Entrada:} A lista \code{vendas} fornecida.
    \item \textbf{Saída:} \code{\{'maçã': 7, 'banana': 7, 'uva': 1, 'laranja': 8\}}
\end{itemize}

\ifmostrarGabarito{
\begin{minted}{python}
# Resposta Questão 6
vendas = ["maçã:5", "banana:3", "maçã:2", "uva:1", "laranja:8", "banana:4"]

consolidado = {}
for item in vendas:
    produto, quantidade = item.split(":")
    consolidado[produto] = consolidado.get(produto, 0) + int(quantidade)

print(consolidado)
\end{minted}
}



% ---------------------------------------------------------
\item Duas palavras são consideradas anagramas se possuem exatamente as mesmas letras nas mesmas quantidades (ex: \code{"amor"} e \code{"roma"}). Crie uma função chamada \code{sao\_anagramas(palavra1,}\\\code{ palavra2)} que retorne \code{True} se as palavras forem anagramas e \code{False} caso contrário.

\textbf{Restrição:} Não utilize métodos ou funções prontas de ordenação (como \code{sorted()}). A validação deve ser feita obrigatoriamente através da comparação de dois dicionários de contagem de caracteres.

Teste a sua função de acordo com as entradas e saídas esperadas a seguir:

\begin{itemize}
    \item \textbf{Entrada:} \code{sao\_anagramas("amor", "roma")} $\rightarrow$ \textbf{Saída:} \code{True}
    \item \textbf{Entrada:} \code{sao\_anagramas("casa", "saco")} $\rightarrow$ \textbf{Saída:} \code{False}
\end{itemize}

\ifmostrarGabarito{
\begin{minted}{python}
# Resposta Questão 7
def criar_freq(palavra):
    freq = {}
    for char in palavra.lower():
        freq[char] = freq.get(char, 0) + 1
    return freq

def sao_anagramas(palavra1, palavra2):
    return criar_freq(palavra1) == criar_freq(palavra2)
\end{minted}
}


% ---------------------------------------------------------
\item Em um sistema administrativo, os dados dos usuários são mantidos em uma lista de dicionários conforme o modelo:
\begin{minted}{python}
usuarios = [
    {"nome": "Carlos", 
    "email": "carlos@email.com", 
    "cpf": "123.456.789-00", 
    "ativo": True},
    {"nome": "Mariana", 
    "email": "mariana@email.com", 
    "cpf": "987.654.321-11", 
    "ativo": False},
    {"nome": "Pedro", 
    "email": "pedro@email.com", 
    "cpf": "456.789.123-22", 
    "ativo": True}
]
\end{minted}
Escreva um código que percorra essa lista e remova a chave \code{"cpf"} do dicionário apenas dos usuários cujo status \code{"ativo"} seja \code{False}, garantindo a privacidade dos dados inativos.
Teste a sua função de acordo com as entradas e saídas esperadas a seguir:

\begin{itemize}
    \item \textbf{Entrada:} A lista \code{usuarios} fornecida.
    \item \textbf{Saída:} 
    \begin{minted}{python}
[
    {'nome': 'Carlos', 
    'email': 'carlos@email.com', 
    'cpf': '123.456.789-00', 
    'ativo': True},
    {'nome': 'Mariana', 
    'email': 'mariana@email.com', 
    'ativo': False},
    {'nome': 'Pedro', 
    'email': 'pedro@email.com', 
    'cpf': '456.789.123-22', 
    'ativo': True}
]
    \end{minted}
\end{itemize}

\ifmostrarGabarito
\begin{minted}{python}
# Resposta Questão 8
usuarios = [
    {"nome": "Carlos", "email": "carlos@email.com", "cpf": "123.456.789-00", "ativo": True},
    {"nome": "Mariana", "email": "mariana@email.com", "cpf": "987.654.321-11", "ativo": False},
    {"nome": "Pedro", "email": "pedro@email.com", "cpf": "456.789.123-22", "ativo": True}
]

for usuario in usuarios:
    if not usuario["ativo"] and "cpf" in usuario:
        del usuario["cpf"]

print(usuarios)
\end{minted}
\fi



% ---------------------------------------------------------
\item Escreva uma função \code{mapear\_posicoes(texto)} que receba um texto e retorne um dicionário onde:
\begin{itemize}
    \item As chaves sejam as palavras únicas presentes no texto (em minúsculas).
    \item Os valores sejam listas de inteiros indicando os índices de posição (iniciando em 0) onde a palavra aparece na sequência das palavras da frase.
\end{itemize}
Teste a sua função de acordo com as entradas e saídas esperadas a seguir:
\begin{itemize}
    \item \textbf{Entrada:} \code{"o sol brilhou e o sol se foi"}
    \item \textbf{Saída:} \code{\{'o': [0, 4], 'sol': [1, 5], 'brilhou': [2], 'e': [3],}\\
    \code{ 'se': [6], 'foi': [7]\}}
\end{itemize}

\ifmostrarGabarito{
\begin{minted}{python}
# Resposta Questão 9
def mapear_posicoes(texto):
    palavras = texto.lower().split()
    mapa = {}
    for idx, palavra in enumerate(palavras):
        if palavra not in mapa:
            mapa[palavra] = []
        mapa[palavra].append(idx)
    return mapa
\end{minted}
}


% ---------------------------------------------------------
\item Crie uma função chamada \code{apurar\_votacao(lista\_votos)} que receba uma lista de strings contendo os nomes dos candidatos votados em uma eleição (podendo conter nomes repetidos ou votos em branco/nulos). A função deve:
\begin{enumerate}[label=\alph*)]
    \item Construir um dicionário com o total de votos recebidos por cada candidato.
    \item Identificar o candidato com o maior número de votos.
    \item Retornar uma tupla contendo o nome do vencedor e o respectivo número de votos no formato \code{(nome\_vencedor, total\_votos)}.
\end{enumerate}

Teste a sua função de acordo com as entradas e saídas esperadas a seguir:
\begin{itemize}
    \item \textbf{Entrada:} \code{["Alice", "Bob", "Alice", "Nulo", "Bob", "Alice", "Branco"]}
    \item \textbf{Saída:} \code{('Alice', 3)}
\end{itemize}

\ifmostrarGabarito{
\begin{minted}{python}
# Resposta Questão 10
def apurar_votacao(lista_votos):
    contagem = {}
    for voto in lista_votos:
        contagem[voto] = contagem.get(voto, 0) + 1
        
    vencedor = max(contagem, key=contagem.get)
    return (vencedor, contagem[vencedor])
\end{minted}
}


\end{enumerate}